%%%%%
%@TheDoctorRAB
%%%%%
%
%%%%%
%standard report format for journal papers
%%%%%
%
%%%%%
%PDF tagging accessible file
%compile with pdflatex
%run $show-pdf-tags --xml filename.pdf
%%%%%
%
%%%%%
%REFERENCES
%neup.bst - numbered citations in order of appearance, short author list with et al in reference section
%nsf.bst - numbered citations in order of appearance, full author list in references section
%standard.bst - citations with author last name with et al for more than 2 authors; full author list in references section
%ans.bst is for ANS only. 
%
%author = {Lastname, Firstname and Lastname, Firstname and Lastname, Firstname} for all bst formats
%bst renders the author list itself
%
%author = {{Nuclear Regulatory Commission}} if the author is an organization, institution, etc., and not people
%
%title = {{}} for all
%
%for all - use \citep{-} - [1] or (Borrelli, 2021) in the text
%standard.bst \cite{-} - Borrelli (2021) in the text
%standard.bst lists references alphabetically
%the rest list numerically
%%%%%


%%%%% enable tagging for accessibility
\RequirePackage{pdfmanagement-testphase}
\DocumentMetadata{
    tagging=on,
    pdfversion=2.0,
    pdfstandard=ua-2,
    lang=en-US,
    debug={log=all},
    testphase={latest,tagpdf},
    tagging-setup={math/setup={mathml-AF,mathml-SE}}
}
\tagpdfsetup{math/alt/use}
%%%%%


%%%%% general 
\documentclass[11pt,a4paper]{article}
\usepackage[lmargin=1in,rmargin=1in,tmargin=1in,bmargin=1in]{geometry}
\usepackage[pagewise]{lineno} %linenumbering restarts each page
%\usepackage[running]{lineno} %linenumbering continuous
\usepackage{mmap} %helps make pdfs copy and paste
\usepackage{setspace}
\usepackage{ulem} %strikethrough - do not \sout{\cite{}}
\usepackage[pdftex,dvipsnames]{xcolor,colortbl} %change font color
\usepackage{graphicx}
\usepackage{tablefootnote}
\usepackage{footnotehyper}
\usepackage{float}
%\usepackage{mypythonhighlight,verbatim}
%\usepackage{subfig}
\usepackage[yyyymmdd]{datetime} %date format
\renewcommand{\dateseparator}{.}
%\graphicspath{{$TEXIMG/}} %path to graphics
\setcounter{secnumdepth}{5} %set subsection to nth level
\usepackage{needspace}
\usepackage[stable,hang,flushmargin]{footmisc} %footnotes in section titles and no indent; standard.bst
\usepackage[inline]{enumitem} %https://www.tug.org/texlive//Contents/live/texmf-dist/doc/latex-dev/latex-lab/latex-lab-enumitem.pdf
\usepackage{boldline}
\usepackage{makecell}
\usepackage{booktabs}
\usepackage{amssymb}
\usepackage{gensymb}
\usepackage{amsmath,nicefrac}
\usepackage{physics}
\usepackage{lscape}
\usepackage{array}
\usepackage{chngcntr}
\usepackage{sectsty}
\usepackage{textcomp}
\usepackage{lastpage}
\usepackage{xargs} %for \newcommandx
\usepackage[colorinlistoftodos,prependcaption,textsize=tiny]{todonotes} %makes colored boxes for commenting
\usepackage{soul}
\usepackage{color}
\usepackage{marginnote}
\usepackage[figure,table]{totalcount}
\usepackage{parskip}
%%%%%


%%%%% links and tags
\usepackage{hyperref}
\hypersetup{
    colorlinks,
    linkcolor=black,
    citecolor=black,
    urlcolor=blue,
    pdftitle={Homework},
    pdfauthor={R. A. Borrelli},
    pdfkeywords={homework},
    pdflanguage={en-US},
    pdfsubject={homework},
    pdfdisplaydoctitle=true 
} 
\usepackage[capitalise]{cleveref} %must be loaded after hyperref
%%%%%


%%%%% tikz
\usepackage{pgf}
\usepackage{tikz} % required for drawing custom shapes
\usetikzlibrary{shapes,arrows,automata,trees}
%%%%%


%%%%% fonts
\usepackage{times}
%arial - uncomment next two lines
%\usepackage{helvet}
%\renewcommand{\familydefault}{\sfdefault}
%%%%%


%%%%% references
%\usepackage[round,semicolon]{natbib} %for (Borrelli 2021; Clooney 2019) - standard.bst 
\usepackage[numbers,sort&compress]{natbib} %for [1-3] - nsf.bst, neup.bst
\setlength{\bibsep}{7pt} %sets space between references
%\renewcommand{\bibsection}{} %suppresses large 'references' heading
%\renewcommand\bibpreamble{\vspace{\baselineskip}} %sets spacing after heading if not using default references heading
%%%%%


%%%%% tables and figures
\usepackage{longtable} %need to put label at top under caption then \\ - use spacing
\usepackage{tablefootnote}
\usepackage{tabularx}
\usepackage{multirow}
\usepackage{tabto} %general tabbed spacing
\usepackage{pdfpages}
\usepackage{wrapfig} %wraps figures around text
\setlength{\intextsep}{0.00mm}
\setlength{\columnsep}{1.00mm}
\usepackage[singlelinecheck=false,labelfont=bf]{caption}
\usepackage{subcaption}
\captionsetup[table]{justification=centering,skip=5pt,labelformat={default},labelsep=period,name={Table}} %sets a space after table caption
\captionsetup[figure]{justification=justified,skip=5pt,labelformat={default},labelsep=period,name={Figure}} %sets space above caption, 'figure' format
\captionsetup[wrapfigure]{justification=centering,aboveskip=0pt,belowskip=0pt,labelformat={default},labelsep=period,name={Fig.}} %sets space above caption, 'figure' format
\captionsetup[wraptable]{justification=centering,aboveskip=0pt,belowskip=0pt,labelformat={default},labelsep=period,name={Table}} %sets space above caption, 'figure' format
%%%%%


%%%%% watermark
%\usepackage[firstpage,vpos=0.63\paperheight]{draftwatermark}
%\SetWatermarkText{\shortstack{DRAFT\\do not distribute}}
%\SetWatermarkScale{0.20}
%%%%%


%%%%% cross referencing files
%\usepackage{xr} %for revisions - will cross reference from one file to here
%\externaldocument{/path/to/auxfilename} %aux file needed
%%%%%


%%%%% toc and glossaries
\usepackage[toc,title]{appendix}
\usepackage[acronym,nomain,nonumberlist]{glossaries}
\makenoidxglossaries
%do not load titlesec,titletoc
%%%%%


%%%%% editing
\newcommand{\edit}[1]{\textcolor{blue}{#1}} %shortcut for changing font color on revised text
\newcommand{\fn}[1]{\footnote{#1}} %shortcut for footnote tag
\newcommand*\sq{\mathbin{\vcenter{\hbox{\rule{.3ex}{.3ex}}}}} %makes a small square as a separator $\sq$
%\newcommand{\sk}[1]{\sout{#1}} %shortcut for default strikethrough - do not sk through citep
\newcommand\sk{\bgroup\markoverwith{\textcolor{red}{\rule[0.5ex]{1pt}{1pt}}}\ULon} %strikethrough with red line; not in \section{}
%\st{} does strikethrough using soul package but does not like acronyms
\newcommand{\blucell}{\cellcolor{aliceblue}} %use to shade in table cell
\newcommand{\grycekk}{\cellcolor{lightgray}} %use to shade in table cell
\newcommand{\whicell}{\cellcolor{antiquewhite}} %use to shade in table cell
%%%%%


%%%%% colors
%http://latexcolor.com/
%https://en.wikibooks.org/wiki/LaTeX/Colors#:~:text=black%2C%20blue%2C%20brown%2C%20cyan,be%20available%20on%20all%20systems.
\definecolor{aliceblue}{rgb}{0.94, 0.97, 1.0}
\definecolor{antiquewhite}{rgb}{0.98, 0.92, 0.84}
\definecolor{lightmauve}{rgb}{0.86, 0.82, 1.0}
\definecolor{brilliantlavender}{rgb}{0.96, 0.73, 1.0}
\definecolor{brandeisblue}{rgb}{0.0, 0.44, 1.0}
\definecolor{darkmidnightblue}{rgb}{0.0, 0.2, 0.4}

\newcommand{\x}{\cellcolor{aliceblue}} %use to shade in table cell
\newcommand{\y}{\cellcolor{lightgray}} %use to shade in table cell
\newcommand{\z}{\cellcolor{antiquewhite}} %use to shade in table cell
%%%%%


%%%%% acronyms
\newcommand{\acf}{\acrfull} %full acronym
\newcommand{\acl}{\acrlong} %long acronym
\newcommand{\acs}{\acrshort} %short acronym

\newcommand{\acfp}{\acrfullpl} %full acronym plural
\newcommand{\aclp}{\acrlongpl} %long acronym plural
\newcommand{\acsp}{\acrshortpl} %short acronym plural
%%%%%


%%%%% todonotes
\newcommandx{\cmt}[2][1=]{\todo[author=\textbf{STRUCTURE},tickmarkheight=0.15cm,linecolor=red,backgroundcolor=red!25,bordercolor=black,#1]{#2}}
\newcommandx{\con}[2][1=]{\todo[author=\textbf{CONTENT},tickmarkheight=0.15cm,linecolor=brilliantlavender,backgroundcolor=brilliantlavender,bordercolor=black,#1]{#2}}
\newcommandx{\rab}[2][1=]{\todo[noline,author=\textbf{RAB},backgroundcolor=Plum!25,bordercolor=black,#1]{#2}}


%\newcommandx{\jon}[2][1=]{\todo[noline,author=\textbf{ATTN: Johnson},backgroundcolor=blue!25,bordercolor=black,#1]{#2}}
%\newcommandx{\han}[2][1=]{\todo[noline,author=\textbf{ATTN: Haney},backgroundcolor=OliveGreen!25,bordercolor=black,#1]{#2}}
%\newcommandx{\rab}[2][1=]{\todo[author=\textbf{RAB},tickmarkheight=0.15cm,linecolor=Plum,backgroundcolor=Plum!25,bordercolor=black,#1]{#2}}
%\newcommandx{\han}[2][1=]{\todo[author=\textbf{ATTN: Haney},tickmarkheight=0.15cm,linecolor=OliveGreen,backgroundcolor=OliveGreen!25,bordercolor=OliveGreen,#1]{#2}}
%\newcommandx{\jon}[2][1=]{\todo[author=\textbf{ATTN: Johnson},tickmarkheight=0.15cm,linecolor=blue,backgroundcolor=blue!25,bordercolor=blue,#1]{#2}}


% highlighting 
\DeclareRobustCommand{\hlc}[1]{{\sethlcolor{LimeGreen}\hl{#1}}}
\makeatletter
    \if@todonotes@disabled
    \newcommand{\hlh}[2]{#1}
    \else
    \newcommand{\hlh}[2]{\han{#2}\hlc{#1}}
    \fi
    \makeatother

\DeclareRobustCommand{\hld}[1]{{\sethlcolor{CornflowerBlue}\hl{#1}}}
\makeatletter
    \if@todonotes@disabled
    \newcommand{\hlj}[2]{#1}
    \else
    \newcommand{\hlj}[2]{\jon{#2}\hld{#1}}
    \fi
    \makeatother

\DeclareRobustCommand{\hlf}[1]{{\sethlcolor{lightmauve}\hl{#1}}}
\makeatletter
    \if@todonotes@disabled
    \newcommand{\hlb}[2]{#1}
    \else
    \newcommand{\hlb}[2]{\rab{#2}\hlf{#1}}
    \fi
    \makeatother
%%%%%


%%%%% table alignments
\newcolumntype{L}[1]{>{\raggedright\let\newline\\\arraybackslash\hspace{0pt}}m{#1}} %uses \raggedright with m,p{} in table column
\newcolumntype{C}[1]{>{\centering\let\newline\\\arraybackslash\hspace{0pt}}m{#1}} %uses \raggedright with m,p{} in table column
\newcolumntype{R}[1]{>{\raggedleft\let\newline\\\arraybackslash\hspace{0pt}}m{#1}} %uses \raggedright with m,p{} in table column
%%%%%


%%%%% table contents
\makeatletter
\renewcommand\tableofcontents{%
    \@starttoc{toc}%
}
\makeatother

\makeatletter
\renewcommand\listoffigures{%
    \@starttoc{lof}%
}
\makeatother

\makeatletter
\renewcommand\listoftables{%
    \@starttoc{lot}%
}
\makeatother

\makeatletter
\newcommand*\ftp{\fontsize{16.5}{17.5}\selectfont}
\makeatother
%%%%%


%%%%% header and footer
\usepackage{fancyhdr}
\pagestyle{fancy}
\fancyhf{} %move page number to bottom right
%\renewcommand{\headrulewidth}{0pt} %turn off line in header
\lhead{\scriptsize NE585}
\chead{\scriptsize \textit{Project 4 -- Nuclear reactor theory}}
\rhead{\scriptsize \today}
\rfoot{\thepage}
%%%%%


%%%%% acronyms
\input{$ACRONYM/acronyms}
%\input{../acronyms/acronyms}
%%%%%


\begin{document}

\begin{titlepage}
    \title{
        NE585 -- Nuclear fuel cycle analysis\\
        Project 4 -- Nuclear reactor theory\\
    }
    \author{
        Name
        \\ \\ \\
        University of Idaho $\sq$ Idaho Falls Center for Higher Education
        \\ \\
        Nuclear Engineering and Industrial Management Department
        \\ \\ \\
        email 
    }
\clearpage %not have page number on title page
\maketitle
\vspace*{\fill}
\begin{flushright}{
        Total -- 400
}
\end{flushright}
\thispagestyle{empty} %start with page number 1 on second page
\end{titlepage}


\section{Neutron multiplication factor -- (25)}
Find the neutron multiplication factor for a fuel of (1) 90\% $^{235}U$ and 10\% $^{238}U$, (2) pure $^{239}Pu$, and (3) pure $^{233}U$ for thermal neutrons. Do the same for fast neutrons ($\sim 2 \; MeV$). Calculate $\eta$; don't just look it up in the book. Assume $\epsilon p$ is approximately unity 


\textit{Comment on any major characteristics of these results in terms of fast v thermal reactors.}





\newpage

\section{Enrichment -- (30)}
What is the minimum $^{235}U$ enrichment needed for a critical assembly of a homogenous core of $^{235}U$ and $^{238}U$? Assume $\epsilon p$ is approximately unity and a thermal core. Comment on the results. 


\textit{This is kind of a trick question.}





\newpage

\section{Fissioning your electricity bill -- (25)}
What is the fission rate, in fissions per day, for a 1000 MWe reactor? Assume 200 MeV is released per fission and 35\% reactor efficiency. If 1.16 atoms of $^{235}U$ are consumed per fission, then how many atoms are consumed per day due to fission? How many grams of $^{235}U$ per day is this? Given how much your last monthly electric bill was, how much $^{235}U$ did you consume per day?  Assume uranium costs (\$100/kg.) 


\textit{Comment on the results.}





\newpage

\section{Neutron flux and current density -- (25)}
Derive the expression for neutron flux and current density for a point source [S n/s] in an infinite vacuum [$D = 0$]. Think physically about what is going on. 





\newpage

\section{Neutron flux and geometry I -- (30)}
If three isotropic point sources of equal strength in an infinite \textit{medium} are located at three corners of an equilateral triangle with sides of length $a$, what is the flux and current at the midpoint of one side of the triangle? The center of the triangle? 


\noindent\textit{D does not equal zero here.}





\newpage

\section{Maximum flux -- (30)}
For the infinite, bare slab geometry of thickness [2a] with source $S \; n/cm^2/s$, \textit{derive} the flux and the maximum flux. In the 2nd edition, it's on p203, eq. 5.41, but should be just around there in the other editions. Neglect the extrapolated distance.





\newpage

\section{Two group theory -- (25)}
For the infinite bare slab geometry again, assume a fast, constant \textit{point source} $s \; n/cm^2/s$, derive explicit expressions for the fast and the thermal flux using two group diffusion theory as follows --
\begin{equation} \label{eq-fast-flux}
    \large D_1\nabla^2{\phi_1}-\Sigma_A^1\phi_1-\Sigma_{1 \rightarrow 2}\phi_1+\Sigma_{2 \rightarrow 1}\phi_2+s=0
\end{equation}

\begin{equation} \label{eq-thermal-flux}
    \large D_2\nabla^2{\phi_2}-\Sigma_A^2\phi_2-\Sigma_{2 \rightarrow 1}\phi_2+\Sigma_{1 \rightarrow 2}\phi_1=0
\end{equation}

\noindent Hints and tips --
\begin{enumerate}
    \item There is only a fast source. Leave it as $s$.
    \item Don't worry about boundary conditions. Leave in the constants and just get the general solution.
    \item Don't assume either absorption cross section is zero.
\end{enumerate}


\textit{This differential equation system is very common to many engineering problems. There is a standard solution approach that can be found with some research. The solution procedure is not tricky, but it can be tedious.}





\newpage

\section{Neutron flux and geometry II -- (30)}
Derive the flux for a rectangular parallelepiped. See the section on the slab reactor, 6.2 in the second edition and assume a separable solution. Show the buckling. Compute maximum to average flux.





\newpage

\section{Critical geometry -- (25)}
What is the critical size of the previous parallelepiped reactor for a homogeneous mixture of $^{239}Pu$ and $Na$, with atom densities of $0.00395 \times 10^{24}$ and $0.0234 \times 10^{24}$, respectively. Make assumptions on the relationship between the x, y, z dimensions. Neglect extrapolation distance. 





\newpage

\section{Critical mass -- (25)}
A spherical thermal reactor of R = 50 cm contains a homogenous core of $^{235}U$ and water. Find the mass of $^{235}U$. It is reflected by graphite. Find the new critical mass. Use room temperature; i.e., standard cross sections.





\newpage

\section{Reactivity equation -- (30)}
Derive the reactivity equation. Work through the derivation in the book (or wherever) and explain the steps sufficiently in terms of the physics of what is going on.
\begin{equation} \label{eq-reactivity}
    \rho=\frac{\omega l_P}{1+\omega l_P}+\frac{\omega}{1+\omega l_P}\cdot\frac{\beta}{\omega+\lambda}
\end{equation}


Prove that $\omega$ is the reciprocal of the stable period with math. Why is there an asymptote at $\rho = 1$? What does that mean physically?





\newpage

\section{Reactor period -- (25)}
Plot reactor period versus absolute value of reactivity for $l_P = 10^{-6},10^{-5},...,10^{-1}$ seconds for the six groups of delayed neutrons for fast fission in $^{235}U$. What observations can be drawn regarding \textit{reactor operation?} Compare to the figure at the beginning of chapter 7 (The Time-Dependent Reactor). It's 7.2 in the 2nd edition.





\newpage

\section{Delayed neutrons -- (30)}
Using the six group reactivity equation for $^{235}U$, show with math that for increasingly negative reactivity, the period is limited (controlled) by the decay of the longest lived precursor. At about what reactivity does this occur? What is the period (\textit{the actual number})?





\newpage

\section{Neutron poison production -- (30)}
Solve for the $^{135}Xe$ and $^{135}I$ concentrations \textit{analytically}, assuming a constant, average thermal flux of $10^{14} n/cm^2/s$. Plot the concentration as a function of time and note any important trends from t = 0 to 500 h. From t = 0 to 200 h, the reactor operates at full power. The reactor is instantaneously scrammed at t = 200 h. Assume a reactor with solely $^{235}U$ with no resonance absorbers or fissionable material. (Just solve the two equations and make whatever simplifying assumptions).





\newpage

\section{Reactor operation -- (15)}
What is the period for a critical reactor?





\newpage

\section{BONUS -- Xenon equilibrium time -- (50)}
What is the equilibrium time for $^{135}Xe$? 

\vspace{\baselineskip}

\textit{Show with math.}





%%%%% header and footer
\pagestyle{fancyplain}
\fancyhf{} %clear header and footer 
\renewcommand{\headrulewidth}{0pt} %set line thickness in header; uncomment as is to remove line
\rfoot{\fancyplain{}{\thepage}}
\nolinenumbers
%%%%%


%%%%%%% citations
%
%%% make sure $BIBREF contains the path to each bib file in .bashrc
%%% make sure Overleaf has the latexmkrc pointing to where the bib and bst are located
%
\bibliographystyle{nsf}
\setlength{\bibhang}{0pt}
\bibliography{
}
%
%%%%%%%


\newpage


\section*{Tables}

{%
\let\oldnumberline\numberline%
\renewcommand{\numberline}{\tablename~\oldnumberline}%
\listoftables%
}


\newpage


%\tagpdfsetup{table/header-rows=1}
%\begin{longtable}{|c|c|}
%        \caption{Average counts for unknown radioactive sample}
%        \label{tab-half-life}
%        \\
%        \hline
%        \textbf{Time (s)}& 
%        \textbf{Counts per second (cps)}
%        \\
%        \hline
%        \endfirsthead
%        0
%        &1000
%        \\
%        30
%        &856.7
%        \\
%        60
%        &734.0
%        \\
%        90
%        &628.8
%        \\
%        120
%        &538.8
%        \\
%        150
%        &461.6
%        \\
%        180
%        &395.4
%        \\
%        210
%        &338.8
%        \\
%        240
%        &290.3
%        \\
%        270
%        &248.7
%        \\
%        300
%        &213.0
%        \\
%        330
%        &182.5
%        \\
%        360
%        &156.4
%        \\
%        390
%        &134.0
%        \\
%        420
%        &114.8
%        \\
%        450
%        &98.3
%        \\
%        480
%        &84.2
%        \\
%        510
%        &72.2
%        \\
%        540
%        &61.8
%        \\
%        570
%        &53.0
%        \\
%        600
%        &45.4
%        \\
%        \hline
%\end{longtable}


\newpage


\section*{Figures}

{%
\let\oldnumberline\numberline%
\renewcommand{\numberline}{\figurename~\oldnumberline}%
\listoffigures%
}


\newpage


%\begin{figure}[h!]
%    \centering
%    \includegraphics[width=0.X\textwidth]{filename}
%    \caption{Descriptive caption}
%    \label{fig-label-name}
%\end{figure}


\begin{appendices}


%%%%% formatting
    \renewcommand{\thesection}{\Roman{section}}
%%%%%


%    \section{\acs{mcp} input decks} \label{app-mcnp}


\newpage


    \section{Acronyms}
    \printnoidxglossary[title={}]


\end{appendices}


\end{document}
